\appendix

\section{Symbols}
\label{app:symbols}

Tables~\ref{tab:symbols} and \ref{tab:symbols2} list the symbols of this paper with the place of
their definition. Symbols of the earlier paper \cite{\SIG{Flat1644}{Flat1671}} keep their meanings. The complete registry, with
the machine checks that no symbol has two meanings, is part of Ref.~\cite{DeflationRepo}.

\begin{table*}
\caption{\label{tab:symbols}The symbols of this paper: geometry and the Dirac field.}
\begin{ruledtabular}
\begin{tabular}{p{0.2\textwidth}p{0.6\textwidth}p{0.14\textwidth}}
symbol & meaning & defined in \\
\hline
$x0,\dots,x7$ & the eight coordinates: $x0$ hidden space-like coordinate; $x1,x2,x3$ ordinary space; $x4$ ordinary time; $x5,x6,x7$ the deflating coordinates, \SIG{time-like}{space-like} & Sec.~\ref{sec:intro} \\
$xt$; $xj$; $a$ & a deflating coordinate, $t\in\{5,6,7\}$; a flat coordinate other than the time, $j\in\{0,1,2,3\}$; index of the four extra axes, $a\in\{0,5,6,7\}$ & Sec.~\ref{sec:intro} \\
$\gff$, $\gss$, $\gvv$; $\rt$ & the three rates of deflation (upright, without subscript; not components of the metric); the rate of $xt$ & Eq.~(\ref{eq:hyp}) \\
$\cff$, $\css$, $\cvv$; $\ct$ & the three positive constants of the metric: the scale factors at $x4=0$ & Eq.~(\ref{eq:hyp}) \\
$\gsum$, $\gsq$, $\gprod$ & the sum of the rates, the sum of their squares, the sum of their products in pairs & Sec.~\ref{sec:geometry} \\
$b_t$ & scale factor of the deflating coordinate $xt$ & Eq.~(\ref{eq:scale}) \\
$\eta$; $g_{\mu\nu}$; $e_A{}^{\mu}$ & flat metric; metric of the deflating spacetime; orthonormal frame & Eq.~(\ref{eq:metric}) \\
$\Vd$ & proper volume of a unit coordinate cell of the deflating coordinates & Eq.~(\ref{eq:Vd}) \\
$\Gamma^{\rho}{}_{\mu\nu}$; $R^{\rho}{}_{\sigma\mu\nu}$, $R^{\mu}{}_{\nu}$, $R$ & Christoffel symbols; Riemann tensor [in Eq.~(\ref{eq:riemann}) its components in the orthonormal frame], Ricci tensor, curvature scalar & Sec.~\ref{sec:curvature} \\
$\cG^{\mu}{}_{\nu}$; $\cH^{\mu}{}_{\nu}$; $\Lam$, $\aGB$ & Einstein tensor; Lanczos tensor of the Gauss-Bonnet term; cosmological constant and coupling of Einstein-Gauss-Bonnet gravity & Sec.~\ref{sec:egb} \\
$\cX_{\mathrm{scale}}$, $\cX_{(s)}$ & the scaling Killing vector; the three special Killing vectors of the case of equal rates & Eqs.~(\ref{eq:scaling}), (\ref{eq:special}) \\
$K_{AB}$; $\bar x4$, $\bar xt$ & sectional curvature of the plane of the frame directions $A$ and $B$; coordinates adapted to the case of equal rates & Appendix~\ref{app:curvature} \\
$P_\mu$; $\varsigma$\SIG{; $x4_{\mathrm{turn}}$}{} & conserved momenta of a classical particle; affine parameter\SIG{; turning value of the ordinary time}{} & Sec.~\ref{sec:geodesics} \\
$\gamma^{A}$, $\gamma^{\mu}$; $\Omega_\mu$, $\omega_{\mu AB}$ & Dirac matrices of the frame and in coordinate components; spin connection and its coefficients & Sec.~\ref{sec:dirac} \\
%%<1644>
$\Cm$; $\chi$, $P^{L,R}$ & the matrix that turns each Dirac matrix into minus its transpose; chirality matrix and the two chiral projectors & Secs.~\ref{sec:dirac}, \ref{sec:clock} \\
%%</1644>
%%<1671>
$\chi$, $P^{L,R}$; $\Cm$ & chirality matrix and the two chiral projectors; the matrix that turns each Dirac matrix into minus its transpose & Secs.~\ref{sec:clock}, \ref{sec:gap} \\
%%</1671>
$\psi_{\tg}$; $\hpsi$ & the 16-component field on the deflating spacetime; the weighted field $\sqrt{\Vd}\,\psi$ & Eqs.~(\ref{eq:dirac}), (\ref{eq:weighted}) \\
%%<1644>
$K$, $K_\perp$; $\gamma^{d}$; $W_{1644}$ & generator of the evolution of a mode in the time, and its part without the scale factors; frame matrix along the wave vector in the deflating coordinates; matrix of the conserved charge & Eq.~(\ref{eq:K}) \\
$\beta$, $\alpha^{j}$; $\bar\psi$ & the Hermitian matrices $-i\gamma^{4}$ and $-\gamma^{4}\gamma^{j}$ of the earlier paper; the adjoint $\psi^\dagger\beta$ & Eq.~(\ref{eq:LVS}) \\
$k_j$, $k_t$; $\Eperp$; $\kd$; $\ud$ & wavenumbers along the flat and the deflating coordinates; $(m^{2}+\sum_jk_j^{2})^{1/2}$; physical wavenumber along the deflating coordinates; $\kd/\rt$ & Eqs.~(\ref{eq:kd}), (\ref{eq:ud}) \\
$w$; $\Swkb$ & one of the eight constant spinors from which the exact solutions are built; exponent of the extinction of a bounded solution & Eqs.~(\ref{eq:exact}), (\ref{eq:Swkb}) \\
%%</1644>
%%<1671>
$\hh$, $\hh_\perp$; $\beta$, $\alpha^{A}$, $\alpha^{d}$; $\bar\psi$ & Hamiltonian of a mode, and its part without the scale factors; $\gamma^{4}$, $\gamma^{4}\gamma^{A}$, and the matrix $\alpha$ along the wave vector in the deflating coordinates; the adjoint $\psi^\dagger\beta$ & Eqs.~(\ref{eq:h}), (\ref{eq:LVS}) \\
$p_j$, $k_t$; $\Eperp$; $\kd$; $\ud$ & wavenumbers along the flat and the deflating coordinates; $(m^{2}+\sum_jp_j^{2})^{1/2}$; physical wavenumber along the deflating coordinates; $\kd/\rt$ & Eqs.~(\ref{eq:kd}), (\ref{eq:ud}) \\
$w$; $Z_+$, $Z_-$ & one of the eight constant spinors from which the exact solutions are built; two cylinder functions whose orders differ by one & Eq.~(\ref{eq:exact}) \\
$\ncr$; $\Tdefl$ & number of quanta that the deflation creates in a mode, per internal state; temperature of their spectrum & Eqs.~(\ref{eq:ncr}), (\ref{eq:Tdefl}) \\
%%</1671>
\end{tabular}
\end{ruledtabular}
\end{table*}

\begin{table*}
\caption{\label{tab:symbols2}The symbols of this paper, continued: the model, the superfluid, and
the cosmological part.}
\begin{ruledtabular}
\begin{tabular}{p{0.2\textwidth}p{0.6\textwidth}p{0.14\textwidth}}
symbol & meaning & defined in \\
\hline
$\mathcal L_{\tg}^{\mathrm{VS}}$; $\PhiVS$, $\hat\Phi_{\mathrm{VS}}$ & Lagrangian of the model; the Vorton-Scalar and its weighted form $\sqrt{\Vd}\,\PhiVS$ & Eq.~(\ref{eq:LVS}) \\
$\lamVS$, $\vVS$, $\yVS$ & self-coupling of the Vorton-Scalar, the constant that fixes its vacuum value, its Yukawa coupling & Sec.~\ref{sec:clock} \\
$U(1)_F$, $\UVS$; $\epsilon$ & the symmetries of fermion number and of chirality; parameter of a symmetry transformation & Sec.~\ref{sec:clock} \\
$m$, $m_s$; $\gVS$, $\aVS$ & fermion mass, mass of the radial mode; coupling of both modes to the fermions, and $\gVS^{2}/4\pi$ & Eq.~(\ref{eq:masses}) \\
$\DVS$, $C$; $D_c$ & the two numbers that decide the state of matter; critical value of $\DVS$ for a bound pair & Eq.~(\ref{eq:DC}) \\
$\td$; $\tpair$ & the deflation clock; the clock measured from the binding of the pair & Eqs.~(\ref{eq:clock}), (\ref{eq:tpair}) \\
$\varphi_{\mathrm{bos}}$, $m_{\mathrm{bos}}$, $\varpi_b$ & a mode of the Vorton-Scalar, its mass, the frequency of its weighted form & Sec.~\ref{sec:scalarmodes} \\
$F_S$, $F_P$; $W(k,k')$ & exchange amplitudes of the radial and of the phase mode; kernel of the gap equation (always written with its two arguments) & Sec.~\ref{sec:gap} \\
$\lambda$, $\lambda_S$, $\lambda_P$; $\Delta_F$; $T_c$ & pairing strength, and its parts from the radial and from the phase mode; gap at the Fermi momentum; critical temperature & Sec.~\ref{sec:gap} \\
$\kF$, $\EF$, $v_F$, $N(0)$ & Fermi momentum, Fermi kinetic energy, Fermi velocity, density of states at the Fermi surface per internal state & Secs.~\ref{sec:gap}, \ref{sec:cjt} \\
$a_{\mathrm{sc}}$, $E_b$ & scattering length of two fermions; binding energy of the pair & Sec.~\ref{sec:pair} \\
$\Xid$; $\Deq$; $\Esh$; $\nex$ & the parameter that decides whether the paired state follows the deflation; equilibrium gap; half-width of the shell of states; density of broken pairs & Sec.~\ref{sec:quench} \\
$\xi_k$, $E_k$, $\kappa$, $n_k$ & single-particle energy, quasi-particle energy, pair correlation and occupation of a pair of states & Eq.~(\ref{eq:bloch}) \\
$V[\Delta]$; $\phi$; $m^{*}$; $T_{1/2}$ & effective potential of the paired state; expectation value of the Vorton-Scalar; fermion mass in matter; temperature at which the fermions have lost half of their mass & Sec.~\ref{sec:cjt} \\
$\Delta^{a}$, $\Delta_t$; $\MM$; $d_\pm$ & order parameter (frame components), and its coordinate components along the deflating coordinates; gap matrix; the two fourfold gaps of a state & Sec.~\ref{sec:order} \\
$SU(2)_A$, $SU(2)_B$, $U(1)_E$\SIG{; $SO(3)_T$}{} & the factors of the internal symmetry of the fields that do not depend on the extra coordinates\SIG{; the rotations of the three additional times $x5,x6,x7$}{} & Sec.~\ref{sec:hidden} \\
$P$, $P_c$; $w_z$; $\theta_{12}$ & generator of $U(1)_E$ and the product of the five matrices of the sector; spinor of a zero mode; angle around a vortex line & Sec.~\ref{sec:vortex} \\
$a_{\mathrm{sp}}$, $H$, $H_0$; $\GN$, $G_8$ & scale factor of ordinary space, Hubble rate, Hubble constant; Newton's constant and the constant of gravitation in eight dimensions & Sec.~\ref{sec:cosmo} \\
$\rho$, $\rho_m$; $\omega_b$, $\omega_c$ & energy density in ordinary space, and that of matter; the densities of baryons and of cold dark matter of the standard model of cosmology & Sec.~\ref{sec:cosmo} \\
$\theta_*$, $r_s$; $t_0$ & angle under which the sound horizon at last scattering is seen, and that sound horizon; age of the universe & Sec.~\ref{sec:cosmo} \\
$\delta_m$, $f_g$, $\sigma_8$ & density contrast of matter, its growth rate, amplitude of clustering & Sec.~\ref{sec:cosmo} \\
$\vEW$, $\nEW$; $\alpha_c$, $\varrho_c$ & electroweak scale and the exponent with which it would follow the coupling under H2; critical coupling of the vacuum condensate, and an abbreviation & Sec.~\ref{sec:higgs} \\
$m_p/m_e$; $G_{\mathrm{Fermi}}$; $\alpha_{\mathrm{em}}$ & ratio of the proton mass to the electron mass; Fermi constant; fine-structure constant & Sec.~\ref{sec:cosmo} \\
H1, H2, H3; A1, A2; HV & the three hypotheses about the superfluid; the two assumptions of the cosmological part; the assumption that the vacuum value of the Vorton-Scalar sets the electroweak scale & Secs.~\ref{sec:intro}, \ref{sec:cosmo} \\
\end{tabular}
\end{ruledtabular}
\end{table*}

\section{Curvature of the deflating metric}
\label{app:curvature}

This appendix derives the Einstein tensor and the Lanczos tensor of Sec.~\ref{sec:geometry} in a
form that can be checked by hand, and gives the coordinates and the additional symmetries of the
case of equal rates.

\subsection{Curvatures on pairs of directions}

In the orthonormal frame the Riemann tensor (\ref{eq:riemann}) of the deflating metric is diagonal
on pairs of directions: for each pair of frame directions $A\ne B$ it has one independent component,
the sectional curvature $K_{AB}=K_{BA}$ of the plane that the two directions span. In the convention
in which space-like components of the metric are positive, which this appendix uses, the sectional
curvatures are, for both signatures,
\begin{equation}
K_{4t}=\rt^{2},\qquad K_{ts}=\rt\,\mathrm{rate}_s\quad(t\ne s),
\label{eq:sectional}
\end{equation}
and $K_{AB}=0$ for every pair that contains one of the flat directions $0,1,2,3$. \SIG{(A sectional
curvature $K_{AB}$ carries two subscripts; it is not the matrix $K$ of Sec.~\ref{sec:dirac}.) }{}For a
Riemann tensor of this kind the mixed components of the Einstein tensor and of the Lanczos tensor
are diagonal, and
\begin{align}
\cG^{A}{}_{A}&=-\sum_{\{B,C\}}K_{BC},\label{eq:pairG}\\
\cH^{A}{}_{A}&=-4\sum_{\{B,C\},\{B',C'\}}K_{BC}\,K_{B'C'},
\label{eq:pairH}
\end{align}
without summation over $A$. The first sum runs over the pairs of directions that do not contain the
direction $A$. The second runs over the unordered couples of two such pairs that have no direction
in common; a set of four directions contributes its three divisions into two pairs. As a test, for a
space of constant curvature in eight dimensions, where all $K_{BC}$ equal one number, the first sum
has $21$ terms and the second $3\times35$, so that the two formulas give $-21$ times that number and
$-420$ times its square, the known values $-\tfrac12\cdot7\cdot6$ and
$-\tfrac12\cdot7\cdot6\cdot5\cdot4$.

\emph{The Einstein tensor.} With Eq.~(\ref{eq:sectional}), Eq.~(\ref{eq:pairG}) gives
\begin{equation}
\begin{aligned}
\cG^{4}{}_{4}&=-\gprod,\qquad\cG^{j}{}_{j}=-(\gsq+\gprod),\\
\cG^{5}{}_{5}&=-(\gss^{2}+\gvv^{2}+\gss\,\gvv),
\end{aligned}
\end{equation}
and the corresponding expressions for $\cG^{6}{}_{6}$ and $\cG^{7}{}_{7}$. Because
$\gsum^{2}=\gsq+2\,\gprod$, these are Eq.~(\ref{eq:einstein}) multiplied by $-\eta_{44}$, the factor
that relates the two conventions.

\emph{The Lanczos tensor.} A couple of pairs contributes to Eq.~(\ref{eq:pairH}) only if neither
pair contains a flat direction, that is, only if the four directions are $4,5,6,7$. For a flat
direction $A=j$ these are allowed, and their three divisions into two pairs give
\begin{equation}
K_{45}K_{67}+K_{46}K_{57}+K_{47}K_{56}=\gff\,\gss\,\gvv\,\gsum ,
\end{equation}
which is Eq.~(\ref{eq:lanczos}). For $A=4,5,6,7$ the direction $A$ is itself among the four, and
these components vanish. The Lovelock tensor of the third order is built in the same way from sets
of six directions, and every such set contains a flat direction: it vanishes.

\emph{The field equations.} In Eq.~(\ref{eq:egb}) the components along $x4$ and along the deflating
coordinates read $\cG^{4}{}_{4}+\Lam=0$ and $\cG^{t}{}_{t}+\Lam=0$. Since
$\cG^{5}{}_{5}-\cG^{6}{}_{6}=(\gff-\gss)\,\gsum$, they can hold together only if the three rates are
equal, and then $\Lam=3\,\gff^{2}$. The components along the flat coordinates read
$-6\,\gff^{2}+\Lam-12\,\aGB\,\gff^{4}=0$, which gives $\aGB=-1/(4\,\gff^{2})$:
Eq.~(\ref{eq:egbsol}). The program verifies Eqs.~(\ref{eq:pairG}) and (\ref{eq:pairH}) for the
deflating metric in exact arithmetic, and an independent program obtains the same Lanczos tensor
from its standard definition by computer algebra, for arbitrary rates and constants.

\subsection{Equal rates: a space of constant curvature}

Let $\gff=\gss=\gvv$, and introduce in place of $x4,\dots,x7$ the coordinates
\begin{equation}
\bar x4=\frac{e^{\gff\,x4}}{\gff},\qquad\bar xt=\ct\,xt\quad(t=5,6,7).
\end{equation}
The part of the line element that contains $x4,\dots,x7$ becomes
\begin{equation}
\eta_{44}\,dx4^{2}-e^{-2\,\gff\,x4}\sum_t d\bar xt^{2}
=\frac{\eta_{44}\,d\bar x4^{2}-\sum_t d\bar xt^{2}}{(\gff\,\bar x4)^{2}} .
\label{eq:conformal}
\end{equation}
\SIG{With $\eta_{44}=-1$ this is minus the line element of four-dimensional hyperbolic space of the
radius $1/\gff$, in its upper-half-space form: the four time-like coordinates span a
negative-definite space of constant curvature.}{With $\eta_{44}=+1$ this is the line element of de
Sitter space of the radius $1/\gff$, in the half of its flat slicing in which space contracts. De
Sitter space is familiar as the geometry of inflation, where the same line element is used with the
opposite direction of time.} The line element (\ref{eq:conformal}) is a multiple of a flat one. The
special conformal transformations of that flat metric along the three axes $\bar x5,\bar x6,\bar x7$
leave the factor $1/\bar x4^{2}$ unchanged and are therefore symmetries. Their generators are, for
$s=5,6,7$,
\begin{equation}
\begin{aligned}
\cX_{(s)}={}&2\,\bar xs\,\Big[\bar x4\,\frac{\partial}{\partial\bar x4}+\sum_t\bar xt\,\frac{\partial}{\partial\bar xt}\Big]\\
&-\Big(\sum_t\bar xt^{2}\SIG{+}{-}\bar x4^{2}\Big)\frac{\partial}{\partial\bar xs},
\end{aligned}
\label{eq:special}
\end{equation}
with $\bar x4\,\partial/\partial\bar x4=\gff^{-1}\partial_4$. Together with the 14 Killing vectors of
the general case and the three rotations of $x5,x6,x7$ they are the 20 Killing vectors of the case
of equal rates.

\section{Machine checks, and how the results are reproduced}
\label{app:checks}

All physics of this paper is computed by one program, written in Rust, which uses the program of the
earlier paper as a library and no other library. Each of its nine commands writes data files and a
list of checks; a check marked ``exact'' is an identity verified in exact rational arithmetic, and a
check marked ``numerical'' states its accuracy. Table~\ref{tab:checks} lists them. The central
results are verified a second time by a program in Python that shares no code with the first and
uses other methods: computer algebra with symbolic rates and constants instead of rational
arithmetic, a general-purpose integrator and the Bessel functions of a library instead of the
integrator of the program. That second program was itself tested by introducing twenty deliberate
errors into it, one at a time; it detected nineteen, and the twentieth, a change of the term
$\gprod$ of Eq.~(\ref{eq:friedmann}), cannot be detected by a comparison of Hubble constants for the
reason given in Sec.~\ref{sec:cosmo} and is covered by the symbolic check of the Einstein tensor.

\begin{table*}
\caption{\label{tab:checks}The machine checks behind this paper: the commands of the program, what
each verifies or computes, and the number of its checks that pass. The checks of both deflating
spacetimes are counted, since the program treats the two together.}
\begin{ruledtabular}
\begin{tabular}{lp{0.68\textwidth}c}
command & content & checks passed \\
\hline
\texttt{matrices} & the Dirac matrices of the frame, verified again & {{GEN:n_flat_matrix_checks}} \\
\texttt{matrices} & geometry, Killing vectors, Einstein and Lanczos tensors, geodesics; matrices in coordinate components, spin connection, Dirac operators & {{GEN:n_deflating_checks}} \\
\texttt{modes} & solutions of the Dirac equation mode by mode: creation ($7+1$) and extinction ($4+4$) & {{GEN:n_mode_checks}} \\
\texttt{symmetry} & surviving rotations and boosts, internal symmetry of the sector, zero modes on a vortex line, reduction of the couplings & {{GEN:n_symmetry_checks}} \\
\texttt{channels} & exchange amplitudes in the pair channels; modes of the Vorton-Scalar & {{GEN:n_channel_checks}} \\
\texttt{history} & gap, bound pair, matter, vacuum and temperature along the clock; order of events & {{GEN:n_history_checks}} \\
\texttt{quench} & dynamics of the paired state under a growing coupling & {{GEN:n_quench_checks}} \\
\texttt{paired} & the paired state and self-bound matter from the effective action & {{GEN:n_paired_checks}} \\
\texttt{cosmology} & inferred Hubble constant, bounds on the rate, the exponent $\nEW$ & {{GEN:n_cosmology_checks}} \\
\texttt{matter} & dark-matter estimates, the resonance, growth of structure, drifts\SIG{}{, temperature of the created quanta} & {{GEN:n_matter_checks}} \\
independent program & geometry, completeness of the Killing vectors, Lanczos tensor, spin connection and Dirac operators by computer algebra; exact solutions; vortex zero modes; Hubble constant & {{GEN:n_deflating_checks_independent}} \\
registry of symbols & no symbol and no name in a data file has two meanings & {{GEN:n_symbol_checks_deflation}} \\
\end{tabular}
\end{ruledtabular}
\end{table*}

No computed number in a table of this paper was typed by hand: the tables are written by a script
from the data files of the program. The figures are drawn by a script from the same files. The
repository \cite{DeflationRepo} contains the programs with their unit tests, the data, the
instructions for compiling and running everything, four supplements and a document on the Dirac
matrices of the deflating spacetimes that derive each result of this paper step by step, and the
record of the errors that were made and corrected during the work.

